We develop a Neural Tangent Kernel (NTK) theory for low-rank and random feature networks (RF‑LR) that provides an approximation principled and computationally 
efficient route to the kernel regime. Assuming fixed weights as classical NTK theory predicts, we prove that low-rank weight matrices
do not loses expressivity. RF‑LR preserves the same reproducing kernel Hilbert space (RKHS) as the shallow ReLU kernel; 
As a toy example, we prove rank‑driven concentration: the empirical three‑layer NTK concentrates around its 
deterministic limit, combining angular and radial decoupling of gaussian processes outputs, we show mean NTK
has the same Laplacian RKHS as the shallow ReLU kernel. At the spectral level for random matrices, under some assumptions to be relaxed, the NTK Gram matrix exhibits a 
spiked and shifted Marchenko–Pastur similar to the 2-layer case : the bulk location is shifted out of 0 thanks to low-rank bottlenecks.
Finally, we give an explicit NTK recursion and closed‑form depth expansion, clarifying how correlation kernels from random features couple across layers.
Taken together, these results establish that RF‑LR preserves expressivity and gives optimization guarantees from smallest eigenvalues
while lowering the entry cost to the kernel regime from \(O(N^2)\) to \(O(rN)\), 
and they yield a tractable framework for finite‑width corrections and spectral predictions at moderate model sizes.
We develop a Neural Tangent Kernel (NTK) theory for low-rank and random feature networks (RF‑LR) that provides an approximation principled and computationally 
efficient route to the kernel regime. Assuming fixed weights as classical NTK theory predicts, we prove that low-rank weight matrices
do not loses expressivity. RF‑LR preserves the same reproducing kernel Hilbert space (RKHS) as the shallow ReLU kernel; 
As a toy example, we prove rank‑driven concentration: the empirical three‑layer NTK concentrates around its 
deterministic limit, combining angular and radial decoupling of gaussian processes outputs, we show mean NTK
has the same Laplacian RKHS as the shallow ReLU kernel. At the spectral level for random matrices, under some assumptions to be relaxed, the NTK Gram matrix exhibits a 
spiked and shifted Marchenko–Pastur similar to the 2-layer case : the bulk location is shifted out of 0 thanks to low-rank bottlenecks.
Finally, we give an explicit NTK recursion and closed‑form depth expansion, clarifying how correlation kernels from random features couple across layers.
Taken together, these results establish that RF‑LR preserves expressivity and gives optimization guarantees from smallest eigenvalues
while lowering the entry cost to the kernel regime from \(O(N^2)\) to \(O(rN)\), 
and they yield a tractable framework for finite‑width corrections and spectral predictions at moderate model sizes.
